\section{Entity and Relation Extraction}
\label{sec:gr-extraction}

Constructing a knowledge graph from unstructured text is fundamentally a
knowledge-acquisition problem: entities and the relations between them must be
identified and normalized before they can be stored as
triples~\cite{ji2021kgsurvey}. This pipeline is typically decomposed into named
entity recognition, relation extraction, and the linking steps that resolve
mentions to canonical graph nodes~\cite{ji2021kgsurvey}. The recent trend is to
drive these steps with large language models, exploiting their broad linguistic
competence to extract structured knowledge with little task-specific
supervision~\cite{pan2024unifying}.

\subsection{Named Entity Recognition}

Named entity recognition (NER) locates spans of text that mention entities and
assigns them a type, and it is the foundation on which relation extraction and
downstream graph population depend~\cite{li2022ner}. Contemporary approaches are
dominated by deep neural architectures that learn contextual representations of
tokens, superseding earlier feature-engineered models~\cite{li2022ner}. In the
Graph RAG setting these extracted entities become the candidate nodes of the
knowledge graph~\cite{edge2024graphrag}.

\subsection{Relation Extraction with LLMs}

Relation extraction determines how the recognized entities are connected,
producing the predicate of each triple. In modern Graph RAG construction this is
frequently performed by prompting a large language model to read source text and
emit entity--relation--entity triples directly, an approach used to derive the
entity knowledge graph in the GraphRAG pipeline~\cite{edge2024graphrag}. Using
language models as extractors is attractive because they generalize across
domains and relation types without dedicated training data, a capability
highlighted in roadmaps for unifying language models and knowledge
graphs~\cite{pan2024unifying}.

\subsection{Coreference Resolution and Entity Linking}

Raw extraction yields many surface mentions that refer to the same real-world
entity, so two normalization steps are required to avoid a fragmented graph.
Coreference resolution groups mentions---including pronouns---that denote the same
entity within and across documents~\cite{lee2017coref}, while entity linking maps
each mention to its canonical identifier in a reference knowledge
base~\cite{shen2015entitylinking}. Performing these steps well is what merges
duplicate nodes into a single coherent entity and gives the resulting graph its
connectivity~\cite{shen2015entitylinking, ji2021kgsurvey}.
